% §1 Introduction --- written last (2026-09-15), for module B's scope: blueprint chapters 8, 9,
% 10, 11 and 14 (ADR-0006); rewritten 2026-09-18 at the author's review, the last section of the
% review, on the same rules as sections 2-8 and on the hub's writing standard. The author's
% finding on the first version: it assumed the line paper had been read in detail and moved very
% fast. So it now starts where its reader is (Gaussian scale space and its axiomatics), says
% what the semigroup axiom excludes and what the hemigroup is, states the line paper's theorem
% with the meaning of its two data, asks the questions a user of the class would ask, and
% answers each in a paragraph with its section; then the literature, the scope, the
% organization, and section 1.1. The separate list of contributions is gone: it repeated the
% five result paragraphs. Nothing here is shared with the blueprint.
%
% Paper I is cited here (one of the three places the reference policy allows, ADR-0003). The line
% paper is cited as its released version. Module C is not cited: it does not exist.
%
% The headline quoted in §1.1 is thm:matern (SpatialLine.matern_theorem): the chapter's own
% headline and the corner the title names; the bridge's strictness (prop:bridge-strictness,
% five declarations on Lean core) is named beside it. The author was asked which to quote and
% the choice is recorded in notes/HANDOFF-article.md for review.
%
% §1.1 owns the trust base of this module (hub writing standard, section 4): what is proved,
% what it rests on (Table 1), what is not formalized, where to look. Its content is unchanged;
% its two long paragraphs are now lists.
%
% Literature: the postdoc's P-0007 round of 2026-09-15 (hub log) --- Becker-Kern's 2004 random
% integral representation is cited for the correspondence with self-similar hemigroups, the
% 2001 max-semistable paper is not needed here, dilative stability is named in §8's open
% questions.

\section{Introduction}\label{sec:introduction}

A scale-space representation replaces a signal by a family of smoothed versions of it, one for
each scale, so that structure can be read at the scale at which it lives. The classical
theory says which smoothing to use. Its requirements are those of a measurement made without
prior knowledge of the signal: linearity, covariance under translations and dilations,
positivity, that a smoothing of a smoothing is again a member of the family, and a requirement
of locality, in one of its forms: that the evolution in scale is generated by a differential
operator, or that no new structure appears as the scale grows. Under them
the kernel must be the Gaussian, and the scale space is the solution of the heat equation
(\sscite{koenderink1984structure}; \sscite{lindeberg1994scale}; \sscite{weickert1999linear}
for the history). The last requirement has a precise form, the \emph{semigroup} property
$K_sK_t = K_{s+t}$ of the smoothing operators $K_t$, and it is stronger than its motivation. It asks that a smoothing of a
smoothing be a member, and it asks in addition that the result depend only on the total amount
of scale crossed. \sscite{pauwels1995extended} showed what the axioms give once the
requirements of locality are set aside: the kernels with transform $e^{-t|\omega|^\alpha}$,
the symmetric stable laws for $0 < \alpha \le 2$. Every one of them other than the Gaussian has
infinite variance.

A companion paper, called the line paper below (\sscite{fagerstrom2026spatial}), kept the
axioms of Pauwels et al.\ and weakened
the semigroup to what its motivation says. The smoothing operators are indexed by pairs of
scales. Call the smoothed signal at a scale the \emph{record} at that scale. The operator
$\Phis{s}{t}$ takes the record at scale $s$ to the record at scale $t$, and the operators are
required to compose,
\[
  \Phis{s}{t}\,\Phis{r}{s} = \Phis{r}{t}, \qquad r \le s \le t ,
\]
a two-parameter \emph{hemigroup}, or cascade. How the stage from $s$ to $t$ depends on the two
scales is left to the other axioms. The main theorem of that paper says what they allow. In
the parametrization of scale in which dilating the signal by $\lambda$ multiplies the scale by
$\lambda$, called canonical below, the kernel at scale $t$ has the Fourier transform
$e^{-F(t\omega)}$, with
\[
  F(\omega) \;=\; a\,\omega^2 + \int_0^\infty \bigl(1 - \cos\omega x\bigr)\,\frac{k(x)}{x}\,dx,
  \qquad a \ge 0, \quad k \ge 0 \ \text{nonincreasing},
\]
and every such $F$ occurs. The two data have a direct meaning. A symmetric kernel is the law
of a random displacement of the point at which the signal is read. The coefficient $a$ is the
Gaussian part of that displacement, and the \emph{profile} $k$ describes its jumps:
displacements of size about $x$ occur at the rate $k(x)\,dx/x$. The Gaussian scale space is
$k = 0$. A profile that is a pure power gives the stable kernels of Pauwels et al. The
exponential profile $k(x) = 2\gamma e^{-x}$ gives the kernels with transform
$(1 + t^2\omega^2)^{-\gamma}$, which for $\gamma > \tfrac12$ are the \Matern{} covariance
functions of spatial statistics. They have every moment finite, and no semigroup that is
covariant under dilation contains them; in the order $\gamma$ at a fixed range they do form
a semigroup, the Bessel scale space discussed below. In probabilistic terms
the admissible kernels are the symmetric \emph{self-decomposable} laws: those of a random
displacement that is, for every contraction ratio, a contracted copy of itself plus an
independent remainder (\sscite{sato1999levy}, Def.~15.1, p.~90). The monotonicity of $k$ is
the whole restriction.

So the class is large, and the line paper stopped at its classification. For an imaging
reader the gain is concrete. The larger class contains kernels with a finite variance and
exponential tails that are computed exactly by cheap recursive filters, which the stable
kernels are not. This paper describes the class, by the questions that a user of it would ask.
What does the class look like as a whole? How does it relate to the time-causal theory, where
the same weakening was carried out first? Which of its members are already known, and what
distinguishes them? What equation does a member solve? And how is a member computed?
\S\ref{sec:line} restates the
line paper's axioms, notation and theorem in the form used here, so that this paper can be
read on its own.

\emph{The shape of the class} (\S\ref{sec:cone}). The admissible exponents $F$ form a convex
cone, and the cone has simple building blocks. The simplest profile is a unit step,
displacements of every size up to a cutoff $\tau$ and none beyond, and its exponent is a
classical special function, the entire cosine integral $\Cin(\tau\omega)$
(\sscite{dlmf2026}, \S6.2). A nonincreasing profile is a stack of unit
steps, so every admissible exponent is a Gaussian term plus a superposition of these, and the
superposition is unique (Proposition~\ref{prop:choquet-cone}, Figure~\ref{fig:admissible-cone}).
The Gaussian and the unit-step members are the extreme rays of the cone, and the named
families are slices through it. The kernel of a unit-step member has no closed form, but it
satisfies a delay equation (Lemma~\ref{lem:cin-delay-equation}, Figure~\ref{fig:cin-rays}).
The equation shows that a hard cutoff in the profile leaves echoes in the kernel at the
multiples of $\tau$. The section ends with a practical criterion. An admissible kernel with no
Gaussian part is bounded at the origin exactly when its profile starts above $1$
(Corollaries~\ref{cor:origin-boundedness} and~\ref{cor:origin-smooth}, the second for a
profile that is unbounded at the origin). For the
\Matern{} family this is the known condition $\gamma > \tfrac12$.

\emph{The relation to the time-causal theory} (\S\ref{sec:bridge}). The same weakening of the
semigroup was first carried out for signals in time, where a measurement may use the past and
not the future (\sscite{fagerstrom2026hemigroup}, called the causal article below). There the
kernels are the laws of random delays on the half-line. The two theories are connected by
Bochner's subordination: run a Brownian motion for a random causal delay, and the displacement
reached has a symmetric law. The section makes this a map from the causal class into the
spatial one, family to family, under which the causal scale is the square of the spatial one
(Proposition~\ref{prop:bridge-families}, Figure~\ref{fig:bridge}). The map is injective, since
the image determines its ancestor, and it is not onto. A kernel in its image has an exponent
whose derivative is positive at every positive frequency, and the exponent of a unit-step
member has stationary points, so the unit-step members have no causal ancestor
(Proposition~\ref{prop:bridge-strictness}). The spatial theory is therefore strictly larger
than the causal theory carried across. The image is the intersection of a decreasing chain of
classes indexed by dimension, the exponents whose radial extension to $\RR^d$ is again
self-decomposable: a member has a causal ancestor exactly when it belongs to every class of
the chain (Proposition~\ref{prop:dimension-filtration}, proved in prose on cited facts about
laws in dimension $d$, which are listed before it). One explicit family shows that the first
two inclusions of the chain are strict (Lemma~\ref{lem:filtration-strictness}, proved here and
not machine-checked).

\emph{The named members} (\S\ref{sec:corners}). Three classical families are located in the
cone (Figure~\ref{fig:corners}). They are called its corners here, in the sense that each is
singled out by an extremal property; they are not extreme rays of the cone, which are the
Gaussian and the unit-step members. The main result concerns the \Matern{} family. Four
descriptions single it out: a profile that is a single exponential, a mixing measure that is a
single atom with no Gaussian part, the transform $(1+t^2\omega^2)^{-\gamma}$ up to the choice of a range, and, for
integer $\gamma$, stages whose transfer function is the $\gamma$-th power of one ratio of
first-order factors (Theorem~\ref{thm:matern}; the family is read in its canonical gauge). Its
members of integer index are exactly the admissible families whose stages, over every finite
ladder of scales, are cascades of identical first-order recursive sections, run forward and
backward. One kernel from the origin of that form already forces the family, and the word
that matters is identical: there are admissible families whose stages are cascades of
first-order sections of different ranges. The \Matern{} members have every moment finite and
exponential tails (Corollary~\ref{cor:matern-realization}). The Student-t laws are the image of
one of the causal families under the bridge, with polynomial tails whose order is set by the
degrees of freedom (Proposition~\ref{prop:student-t}). The symmetric stable laws are the
members for which the stages depend only on the difference of their scales, and the heavy
tails of all of them but the Gaussian are the price of that requirement
(Proposition~\ref{prop:stable-family}). Two general
facts make the comparison possible. The number of finite moments of a kernel and the weight
of its tail are read off the tail of the profile (Proposition~\ref{prop:moments-tails}). And
all three families lie in one subcone, the profiles that are mixtures of exponentials, on
which the bridge is a bijection (Proposition~\ref{prop:thorin-subclass}). For the Student-t
family the membership rests on a theorem of Halgreen's that is taken as a hypothesis.

\emph{The evolution equation} (\S\ref{sec:generator}). For a Schwartz signal, every member's
scale space solves an equation $\partial_t u = \mathcal{A}_t u$. The generator is a Laplacian
for the Gaussian coefficient plus, for every displacement size in the catalogue, the amount by
which the average of the signal's two neighbours at that distance exceeds the signal
(Proposition~\ref{prop:scale-evolution}, Figure~\ref{fig:generator}). The heat equation and
the fractional heat equation of the $\alpha$-scale spaces (\sscite{duits2004axioms}) are
recovered for the Gaussian and the stable family. The \Matern{} family has a bounded generator, and a unit-step member has a
single second difference (Proposition~\ref{prop:corner-generators}).

\emph{The algorithm} (\S\ref{sec:implementation}). The hemigroup is also how a member is
computed. Over any finite ladder of scales, the record at each scale is obtained from the
record at the previous one by one stage, and the result is the scale space itself, with no
approximation in the scale direction (Proposition~\ref{prop:knot-exactness},
Figure~\ref{fig:sections}). For the \Matern{} family with integer $\gamma$ a stage is $\gamma$
pairs of first-order recursive filters, each pair with one filter run forward along the signal
and one backward (Proposition~\ref{prop:matern-cascade}). This is the kind of filter long used
to approximate the Gaussian (\sscite{deriche1993recursively,young1995recursive}), and the
theory says which of those filters are themselves scale spaces
(Proposition~\ref{prop:rational-criterion}, read at the two designs in
Remark~\ref{rem:recursive-gaussians}). The members whose profile is a mixture of exponentials,
the stable family among them and the Student-t family under the hypothesis named above, are
limits of members with finitely many ranges, and a member with finitely many ranges and
integer weights has such stages (Proposition~\ref{prop:thorin-machine}). The families realized
by such sections reach, in the limit, exactly the members whose Thorin measure consists of
atoms of even integer mass, finitely many in every bounded set, possibly with a Gaussian part.
Every such limit has an exponential moment, so the stable and the Student-t members, which
lack moments, are not among the limits, and this needs no hypothesis
(Proposition~\ref{prop:rational-closure}). The kernels from the origin of the limits are the
symmetric P\'olya frequency functions of Schoenberg's theorem
(Remark~\ref{rem:polya-limits}). A numerical run closes the section
(Example~\ref{ex:numerical-run}). It checks the statements on a sampled signal, it measures
how far the members with such stages stay from a member like the Poisson scale space, and it
finds that neither of the two recursive Gaussians in common use is the kernel of an exact
cascade.

\emph{Relation to the literature.} The families at the corners are classical, and the
section on each says where. The \Matern{} kernels are the covariance functions of
\sscite{whittle1954stationary,guttorp2006studies,porcu2024matern}. As a scale-space family
composed along the smoothness and not along the range they are the Bessel scale space of
\sscite{burgeth2005bessel}, who also dilate the kernels at a fixed order; what the hemigroup
adds there is the stage between two ranges, which makes the family in the range a cascade. The
same authors' relativistic scale spaces (\sscite{burgeth2005relativistic}), with the
exponent $\sqrt{\omega^2 + m^2} - m$, are a further non-Gaussian semigroup in imaging with
exponentially decaying kernels; their scale parameter is not the dilation of the kernel, so
they are covariant in a different sense from the one required here. In probability the
\Matern{} kernels are the variance-gamma laws of
\sscite{madan1990variance,fischer2025variance}. The mixtures of exponential profiles are the
symmetric part of the generalized gamma convolutions of Thorin and Bondesson
(\sscite{bondesson1992generalized}), and the Student-t corner rests on Halgreen's theorem
that the generalized inverse Gaussian laws belong to that class (\sscite{halgreen1979self}).
The subordination bridge is Bochner's, and its Bernstein-function calculus is that of
\sscite{schilling2012bernstein}, where Schoenberg's radial theorem, used for the chain of
classes by dimension, is also stated. The Gaussian variance mixtures with an infinitely
divisible mixing law, the type G laws, are the class the image of the bridge sits inside
(\sscite{maejima2002type}; \sscite{steutel2004infinite}). Self-decomposable laws and self-similar additive processes are
the probability classes of \sscite{sato1999levy,sato1991selfsimilar}. Hemigroups appear in the
probability-on-groups literature
(\sscite{heyer1977probability,siebert1982continuous,hazod2001stable}), and the random
integral representation of a self-decomposable law
(\sscite{jurek1983integral}; \sscite{beckerkern2004random} for hemigroups) is the closest in
shape to the bridge, though it starts from no axioms on measurement. In the scale-space
literature the nearest construction is Lindeberg's time-causal limit kernel
(\sscite{lindeberg2016time}), an infinite cascade of first-order filters over a geometric
ladder of time constants, of which the pyramid of \S\ref{sec:rational} is the two-sided form. The $\alpha$-scale spaces of \sscite{duits2004axioms} are the evolution
equation of the stable corner, and the Poisson scale space of \sscite{felsberg2004monogenic}
is its member of index $1$.

What this paper adds is the geometry of the class as a consequence of the axioms, and the
bridge as a map between two theories with its strictness proved. It adds the
characterization of the \Matern{} family by its exact implementation, the evolution equation
of every member, and the implementation as a theorem about the cascade.

\emph{Scope.} Locality of the evolution and the non-creation of local extrema across scale,
the requirements from which the Gaussian axiomatics select one member, are not treated here.
They are the subject of the next module, and nothing below depends on them. The paper stays on
the line and in the $L^1$ signal model, and it makes no empirical claim about natural signals.
Most of the results of \S\S\ref{sec:cone}--\ref{sec:generator} are machine-checked in the
Lean~4 proof assistant, and the results of \S\ref{sec:implementation} are proved in prose.
\S\ref{sec:machine-checked} says which, and what they rest on. The remaining sections mention
the formalization only where a statement or a step is cited and not proved.
\S\ref{sec:conclusions} collects what is settled and what is left open.

\subsection{What is machine-checked}\label{sec:machine-checked}

Most of the results of \S\S\ref{sec:cone}--\ref{sec:generator} have been formalized in the
Lean~4 proof assistant, on top of its mathematical library Mathlib, in the same development as
the line paper's. The exceptions are listed below. This subsection is self-contained, and
a reader interested in the mathematics alone can pass to \S\ref{sec:line}.

\paragraph{What is proved.} Section by section:
\begin{itemize}
\item of \S\ref{sec:cone}, every proposition, lemma and corollary but the last: the $\Cin$
rays and the superposition, the Choquet structure in all its parts, the delay equation in both
forms, the folding translation, and the behaviour at the origin with the boundedness threshold
for a profile that is finite there;
\item of \S\ref{sec:bridge}, everything but the dimension filtration and its lemma: the bridge
on exponents in both clauses, the bridge on families in its five clauses, that of the Student-t
family being the mixture identity, with the self-decomposability of its delay law as a
hypothesis, the
strictness proposition in all four clauses, and the memberships of the Gaussian ray and the
\Matern{}, Student-t and stable families, the last under its hypothesis;
\item of \S\ref{sec:corners}, the \Matern{} theorem in its four clauses, the \Matern{}
exponent, transforms and moments, the Thorin subclass in its representation, its bridge and
its strictness, the two-sided reading of Thorin's exponent, the Student-t family in its
kernel, its transform and its Thorin clause under its hypothesis, and the symmetric stable
family in all its parts; and of the moments-and-tails proposition, the variance formula, the
completely monotone case and the Gaussian exclusion;
\item of \S\ref{sec:generator}, everything. The generator of the Thorin members uses the
uniqueness of the \Levy--Khintchine pair, through the representation of \S\ref{sec:thorin};
the rest of the section rests on Lean's logical foundations alone.
\end{itemize}
The headline theorem is assembled in one declaration, so that it is verified as it is printed.

\paragraph{What it rests on.} Beyond Lean's logical foundations, the results above use the
cited facts of the line paper and three more. The line paper's are Bochner's theorem in the
forward direction, the symmetric \Levy--Khintchine representation at its converse and
uniqueness clauses, and the regularity of a self-decomposable density. The three further
facts are each cited with a page reference checked against the source where they are used.
\begin{enumerate}
\item \emph{Bernstein's theorem, the existence equivalence only.} A function on the half-line
is completely monotone if and only if it is the Laplace transform of a positive measure
(\sscite{schilling2012bernstein}, Thm.~1.4, p.~3). The uniqueness of the measure, which the
source states with it, is proved and not cited (Proposition~\ref{prop:laplace-uniqueness-locally-finite}).
\item \emph{The moment criterion for infinitely divisible laws and the tail floor for a
bounded jump support} (\sscite{sato1999levy}, Thm.~25.3, p.~159, and Thm.~26.1, p.~168), the
second at its divergence branch alone and at a lower bound for the radius.
\item \emph{The behaviour of a self-decomposable density at the origin} (\sscite{sato1999levy},
Thm.~28.4, p.~191, and Thm.~53.8, pp.~410--411), three regimes admitted as three named facts
at the source's own letter, over a two-sided profile and its two one-sided limits.
\end{enumerate}
Further cited facts are used by statements that are proved in prose and not machine-checked.
Proposition~\ref{prop:dimension-filtration} uses Schoenberg's radial theorem
(\sscite{schilling2012bernstein}, Thm.~13.14, p.~212), in two of its clauses, with the
\Levy--Khintchine representation and the polar criterion of self-decomposability in dimension
$d$ and the cited facts listed before it, and the third clause of
Lemma~\ref{lem:filtration-strictness} uses the same two facts in dimension $d$.
Proposition~\ref{prop:rational-closure} uses Bondesson's closure theorem for the generalized
gamma convolutions (\sscite{bondesson1992generalized}, Thm.~3.1.5, pp.~34--35), with the
continuity theorem for Laplace transforms (\sscite{feller2009introduction}, Vol.~2, \S XIII.1,
Thm.~2, p.~431), and the moment criterion a second time.
Corollary~\ref{cor:origin-smooth} uses a further clause of the theorem on the density at the
origin (\sscite{sato1999levy}, Thm.~28.4(ii), p.~191), which the development admits nowhere.
Table~\ref{tab:trust-base} lists every cited fact with the statements that use it directly.

\begin{table}[!htb]
\centering
\caption{The cited facts and their direct uses, with the entry of the development's ledger
that carries each. A numbered proposition, lemma, theorem or corollary that is not listed uses
no cited fact directly, apart from the regularity of a self-decomposable density, which enters
through Proposition~\ref{prop:kernel-regularity}. Cited and used by no machine-checked proof:
the closed form of the \Matern{} kernel and the Bessel asymptotics behind its tail (ledger
entries A15 and A16, Proposition~\ref{prop:matern-density}, from which
Corollary~\ref{cor:matern-realization} reads its exponential tails); Thorin's definition of
the extended generalized gamma convolutions (A17, Lemma~\ref{lem:thorin-ggc}); Halgreen's
theorem on the inverse-gamma law (A18, Remark~\ref{rem:student-ggc}). Cited in remarks, for
orientation or attribution: the identification of the constant of Lemma~\ref{lem:cin-rays}
(\sscite{dlmf2026}, Remark~\ref{rem:cin-constant}), the background driving process
(\sscite{sato1999levy}, Thm.~17.5), and Schoenberg's theorem on P\'olya frequency functions
(A21, Remark~\ref{rem:polya-limits}). The rows below the rule are used only by statements
proved in prose.}
\label{tab:trust-base}
\footnotesize
\setlength{\tabcolsep}{3pt}
\begin{tabular}{@{}lllp{4.9cm}@{}}
\toprule
cited fact & entry & clause used & used directly by \\
\midrule
Bochner's theorem & A1 & forward direction & Prop.~\ref{prop:bridge-families} \\
symmetric \Levy--Khintchine & A3 & converse & Prop.~\ref{prop:bridge-families} \\
representation & & uniqueness & Prop.~\ref{prop:choquet-cone}; Thm.~\ref{thm:matern}; Prop.~\ref{prop:thorin-subclass}; \\
 & & & Prop.~\ref{prop:student-t}(3); Prop.~\ref{prop:corner-generators} \\
Bernstein's theorem & A11 & existence equivalence & Prop.~\ref{prop:thorin-subclass}; Prop.~\ref{prop:moments-tails}(2), last sentence; \\
 & & & Prop.~\ref{prop:student-t}(3) \\
the moment criterion & A13 & as stated & Prop.~\ref{prop:moments-tails}(1), which is the fact; Prop.~\ref{prop:stable-family} \\
the tail floor & A14 & divergence, reached radius & Prop.~\ref{prop:moments-tails}(2), last sentence \\
the density at the origin & A10 & three regimes, at the letter & Prop.~\ref{prop:cin-origin-singularity}; Cor.~\ref{cor:origin-boundedness} \\
\midrule
the density at the origin & A10 & the clause at $k(0{+}) = \infty$ & Cor.~\ref{cor:origin-smooth} \\
the moment criterion & A13 & exponential weight & Prop.~\ref{prop:rational-closure}(5) \\
\Levy--Khintchine in dimension $d$ & A3 & converse, uniqueness & Prop.~\ref{prop:dimension-filtration}; Lem.~\ref{lem:filtration-strictness}(3) \\
the polar criterion in dimension $d$ & A7 & both directions & Prop.~\ref{prop:dimension-filtration}; Lem.~\ref{lem:filtration-strictness}(3) \\
the radial theorem & A23 & Bernstein from every dimension & Prop.~\ref{prop:dimension-filtration}(3), (4) \\
closure of the gamma convolutions, & A24 & closure, vague convergence & Prop.~\ref{prop:rational-closure}(1), (4) \\
with continuity of Laplace transforms & & & \\
\bottomrule
\end{tabular}
\end{table}
The three names of the fact on the density at the origin carry the same steps beyond their
pages. The profile is any monotone two-sided profile and not the source's version of it that
is continuous from one side, which differs from it on a countable set and gives the same law.
The hypothesis is this paper's cosine-transform specification of a symmetric law and
not the source's self-decomposable law, a passage through the \Levy--Khintchine
representation already in the trust base. The conclusion is a two-sided comparison on a
punctured neighbourhood with an existential representative, against the source's comparison
function with its slowly varying factor. That is less than the source's asymptotic with its
constant because the constant is not zero, which holds for a symmetric law. The passage between the source's two-sided convention and this paper's folded one
is a theorem beside them (Lemma~\ref{lem:folding-translation}). An earlier version of the
development had admitted the singular regime without the slowly varying factor, which made
the admitted fact false; an external review found it, and the fact is now stated as the source
states it (\sscite{sato1999levy}, (53.25) and (53.28), p.~411). In print the repair added the
slowly varying factor and the comparison function of the threshold regime to the statement of
Proposition~\ref{prop:cin-origin-singularity}. The threshold of
Corollary~\ref{cor:origin-boundedness} and every later result were unchanged.

\paragraph{What is not formalized.} The following, in the scope of this paper, and the text
names each where it occurs.
\begin{itemize}
\item The smoothness of the kernel when the profile is unbounded at the origin, proved in
prose on a clause of a cited theorem (Corollary~\ref{cor:origin-smooth}).
\item The moment criterion at bounded jump support and the finiteness branch of the tail
floor (Proposition~\ref{prop:moments-tails}, two clauses).
\item The closed form of the \Matern{} kernel and its tail
(Proposition~\ref{prop:matern-density}).
\item The Thorin subclass as Thorin's own class, which is the reading of a definition
(Lemma~\ref{lem:thorin-ggc}).
\item The moment count of the Student-t family (Proposition~\ref{prop:student-t}(4)).
\item The dimension filtration, a proposition proved in prose on cited facts about laws in
dimension $d$, of which this development has none
(Proposition~\ref{prop:dimension-filtration}).
\item Its lemma, of which two clauses are stated in the development and not proved there, and
the third is proved in prose (Lemma~\ref{lem:filtration-strictness}).
\item The whole of the implementation section, \S\ref{sec:implementation}, where no
machine-checked proof is scheduled, as for the causal article's, and the corollary of the
\Matern{} theorem that rests on it (Corollary~\ref{cor:matern-realization}). Its statements are
proved in prose from the statements of this paper, except
Proposition~\ref{prop:rational-closure}, which is proved on cited facts.
\end{itemize}
Every other numbered proposition, lemma, theorem and corollary is machine-checked. Who has read
the statements of this list is said in the statement on the use of AI at the end of the paper.
The remarks are commentary. One of them describes a computation that is not
carried out (Remark~\ref{rem:cin-propagation}), and two carry the cited facts that ground
hypotheses and not statements (Remarks~\ref{rem:bridge-subordination}
and~\ref{rem:student-ggc}).

\paragraph{Where to look.} The development and its axiom ledger, with page-verified citations
for every cited fact, are at the repository accompanying this paper,
\url{https://github.com/danielfagerstrom/spatial-hemigroup-scale-space-cone}, which is deposited
with this version of the paper under the DOI printed on the first page. It is an export from
the development repository restricted to what the statements of this paper rest on: the Lean
modules, the trust boundary and the ledger, and the sources of the blueprint chapters behind
\S\ref{sec:cone}--\S\ref{sec:implementation}, in which the declaration behind each numbered
statement is named in a \texttt{\char`\\lean} tag on that statement. The toolchain is pinned
in-tree, so the verification is reproducible with one command, \texttt{lake build} in
\texttt{Formalization/}. A second, \texttt{lake env lean CIAxiomGuard.lean}, prints the axiom
usage of every named declaration, to be read against the trust-boundary file beside it. On the
declaration \texttt{SpatialLine.matern\_theorem}, the headline of \S\ref{sec:corners}, it prints
Lean's three logical axioms and the uniqueness clause of the \Levy--Khintchine representation,
and nothing else:
\begin{verbatim}
'SpatialLine.matern_theorem' depends on axioms: [propext,
 Classical.choice, Quot.sound,
 SpatialLine.fourier_toolbox_levy_unique]
\end{verbatim}
The strictness of the bridge, Proposition~\ref{prop:bridge-strictness}, prints the three
logical axioms alone in all five of its declarations.
